\section{Introduction}\label{sec:intro}

A logarithmic-space machine can remember a bounded number of positions
in its input, but cannot keep a polynomial-length computation or random
tape. Can fresh random bits increase what such a machine can decide?
We answer this question by giving a deterministic logarithmic-space
simulation of every polynomial-time randomized logarithmic-space
computation with bounded error.

We use machines with finite control, an endmarked read-only input,
finitely many work tapes, and fresh independent fair bits. Input heads
are confined between the endmarkers. All randomized running-time
bounds hold on every random tape. Space counts traversed writable cells
in bits, including every simultaneously live counter, numerical register
and suspended call; visiting a blank work cell counts toward space.
A transducer has a one-way write-only output tape. Throughout,
$n$ is the input length and logarithms have base two.

The class $\Lclass$ consists of languages decidable by deterministic
machines in $O(\log(n+2))$ work space. Such a halting deterministic
machine has only polynomially many configurations and hence runs in
polynomial time. The class $\RL$ permits
polynomial-time randomized machines in this space, with acceptance
probability zero on no inputs and at least $1/2$ on yes inputs.
For $\BPL$, the corresponding probabilities are at most $1/3$ and at
least $2/3$.

\begin{theorem}[Exact logarithmic-space derandomization]\label{thm:main}
\[
                 \Lclass=\RL=\BPL.
\]
\end{theorem}

The proof approximates arbitrary acceptance probabilities. This gives
the two-sided statement directly, and also applies to computations
whose acceptance probabilities have no decision gap. Its quantitative
form separates the memory needed for the input from the memory needed
for the requested accuracy.

\begin{theorem}[Acceptance-probability approximation]\label{thm:probability}
Fix a randomized machine $\mathcal M$ with polynomial worst-case
running time and $O(\log(n+2))$ work space. Let $p_{\mathcal M}(x)$ be
its acceptance probability on an input $x$ of length $n$. There is a
uniform deterministic algorithm which, on input $(x,1^q)$ with $q\ge1$,
outputs an integer $0\le a\le2^{q+2}$ satisfying
\[
       \left|a2^{-(q+2)}-p_{\mathcal M}(x)\right|\le2^{-q}.
\]
It uses $O(\log(n+2)+q)$ work space and
$(n+2)^{O(1)}2^{O(q)}$ time. The constants depend only on $\mathcal M$.
\end{theorem}

In particular, every fixed inverse-polynomial accuracy is attained in
logarithmic space and polynomial time. The precision parameter is unary;
the running-time bound is exponential in general $q$. We first prove
the fixed inverse-polynomial result and then obtain the single
variable-precision algorithm by padding. Section~\ref{sec:conclusion}
also derives total deterministic separators for bounded-error promise
problems and constructs accepting runs when the machine's acceptance
probability has a specified inverse-polynomial lower bound.

The simulation is also effective at the level of machine descriptions.
Given a source program and supplied polynomial-time and logarithmic-space
bounds, one terminating compiler produces a deterministic decider and
explicit numerical bounds for its space and time
(Theorem~\ref{thm:effective-simulation}). The supplied bounds are
hypotheses on the source, rather than properties the compiler must decide.
Section~\ref{sec:compiler} constructs this compiler by specializing one
fixed deterministic library.

\subsection{Historical context}

The complexity of probabilistic Turing machines was developed by
Gill~\cite{Gill1977}, including measures of both time and tape usage.
Random walks made the logarithmic-space question concrete:
Aleliunas, Karp, Lipton, Lov\'asz and Rackoff obtained polynomial-time
randomized algorithms for undirected connectivity and explicitly asked
whether polynomial-time randomized logarithmic space could be simulated
in deterministic logarithmic space~\cite{AleliunasKLLR1979}.
Early work also studied probabilistic machines without a polynomial-time
restriction. Borodin, Cook and Pippenger obtained a quadratic-space
deterministic simulation in this broader setting
\cite{BorodinCookPippenger1983}.
The distinction between a space bound and simultaneous time and space
bounds is essential throughout this history.

Pseudorandom generators attack the problem by replacing a long sequence
of fresh random bits with the output generated from a short seed.
Ajtai, Koml\'os and Szemer\'edi gave early logarithmic-space
constructions that hit small-space tests using short random
strings~\cite{AjtaiKomlosSzemeredi1987}. Babai, Nisan and Szegedy
subsequently obtained pseudorandom generators through multiparty
communication bounds~\cite{BabaiNisanSzegedy1992}.
Nisan's generator uses $O(S\log R)$ seed bits for a space-$S$
computation reading $R$ random bits~\cite{Nisan1992}.
For logarithmic space and polynomially many bits this gives
$O(\log^2 n)$ seed length; direct seed enumeration has a
quasipolynomial time bound. Nisan's later simultaneous simulation
achieves polynomial time and $O(\log^2 n)$ space
\cite{Nisan1994RLinSC}. Although its title uses $\mathrm{RL}$,
that result includes two-sided bounded error.
For $S(n)\ge\log n$, Nisan and Zuckerman showed that a space-$S$
computation using $\operatorname{poly}(S)$ random bits can be simulated
using only $O(S)$ random bits and $O(S)$ space~\cite{NisanZuckerman1996}.
The restriction on the original number of random bits is substantive;
for logarithmic space it permits polylogarithmically many bits.

Saks and Zhou reduced the deterministic space needed for general
bounded-error logarithmic-space computation to
$O(\log^{3/2}n)$~\cite{SaksZhou1999}.
Cai, Chakaravarthy and van Melkebeek established a simultaneous
tradeoff: for each fixed rational $0\le\alpha\le1/2$, the simulation
uses time $n^{O(\log^{1/2-\alpha}n)}$ and space
$O(\log^{3/2+\alpha}n)$
\cite{CaiChakaravarthyVanMelkebeek2006}.
The polynomial-time endpoint has squared-logarithmic space.
Hoza subsequently improved the general space bound to
$O(\log^{3/2}n/\sqrt{\log\log n})$
\cite[Theorem~1.6 and Corollary~A.5]{Hoza2021}.
These smaller space containments do not assert the same polynomial-time
endpoint. Hoza's work also develops weighted pseudorandomness, in which
signed or real weights replace an ordinary probability distribution;
the space-containment improvement and the new weighted-generator
construction are distinct results of that paper.

Acceptance-probability estimation also connects derandomization to
stochastic matrix products and inverses. Ahmadinejad, Kelner, Murtagh,
Peebles, Sidford and Vadhan developed high-precision small-space
algorithms for random-walk probabilities, with stronger bounds for
Eulerian graphs~\cite{AhmadinejadKelnerMurtaghPeeblesSidfordVadhan2020}.
Cohen, Doron, Sberlo and Ta-Shma later used varying precision across
recursion to improve the space complexity of long stochastic matrix
products~\cite{CohenDoronSberloTaShma2023}. The construction below
works with a finite inverse on a time-layered configuration graph;
its correction and copy hierarchy controls the information retained
while individual entries are evaluated.

Structural restrictions on a computation provide another route.
For undirected connectivity, Nisan, Szemer\'edi and Wigderson obtained
deterministic space $O(\log^{3/2}n)$~\cite{NisanSzemerediWigderson1992},
subsequently improved to $O(\log^{4/3}n)$ by Armoni, Ta-Shma,
Wigderson and Zhou~\cite{ArmoniTaShmaWigdersonZhou2000}.
Reingold's logarithmic-space algorithm, whose conference version
appeared at STOC 2005, settles undirected connectivity~\cite{Reingold2005,Reingold2008}.
Reingold, Trevisan and Vadhan studied regular directed graphs and gave
logarithmic-space path search for Eulerian directed graphs
\cite{RTV2006}. Their oblivious walk constructions impose additional
labeling conditions. General directed configuration graphs need not
have these properties, so those results did not settle the unrestricted
randomized computation problem.
A read-once branching program is a layered deterministic computation
that reads the input symbols once in a fixed order (bits in the binary
case). Its length is the number of steps,
and its width is the largest number of states in a layer.
A $1/2$-hitting set for a class of programs contains an accepting input
for every member that accepts at least half of all input strings.
The relation between one-sided and two-sided derandomization also
requires care. Cheng and Hoza proved that one uniformly logarithmic-space-enumerable
$1/2$-hitting-set family for all width-$n$, length-$n$ read-once branching
programs would yield
$\mathsf L=\mathsf{BPL}$~\cite{ChengHoza2022}.
The premise is a hitting-set construction; it is not merely the
language-class equality $\mathsf{RL}=\mathsf L$.

Recent results refine both general and structured models.
Chen, Cohen, Doron, Khaskelberg and Ta-Shma improve error reduction for
weighted generators; in the polynomial-width regime their construction
has seed length $\log|\Sigma|+O(\log^2 T+\log(1/\delta))$, where
$T$ is branching-program length, $\Sigma$ its alphabet and $\delta$
the error~\cite{ChenCohenDoronKhaskelbergTaShma2026}.
Cheng and Wu obtain simultaneous polynomial-time and polylogarithmic-space
simulations for regular branching programs, with extensions to polynomial-time
randomized logarithmic-space machines whose two-way read-only random tape
visits each cell at most constantly often, and to such machines with read-once randomness
and an auxiliary stack whose push, pop and idle schedule is independent
of the random bits. Consecutive stays on a random-tape cell count as one
visit~\cite[Theorems~5--6]{ChengWu2026SC}.
Hardness-versus-randomness approaches additionally give conditional,
certified and per-instance alternatives; Pyne and Tell survey this
line of work~\cite{PyneTell2026}.

Effectivity has its own substantial antecedents.
Pyne, Raz and Zhan construct a universal deterministic estimator whose
$O(S(n))$ space guarantee follows from
$\mathsf{prBPL}\subseteq\mathsf{SPACE}[O(S(n))]$, for
space-constructible $S(n)\ge\log n$
\cite[full version, Theorem~5.1]{PyneRazZhan2023}.
Here $\mathsf{prBPL}$ denotes promise problems solvable in polynomial
time and randomized logarithmic space with bounded error.
The compiler below uses a fixed separator established in this paper
and supplies explicit source-specific numerical resource bounds.
Its construction should be read in this context of universal
simulation, with the source bounds supplied as part of the compiler input.

\subsection{Proof strategy}

Fix a randomized machine and an input of length $n$, and write
$B=\Theta(\log(n+2))$.  The proof constructs a finite family of estimates of its
acceptance probability, indexed by $O(B)$-bit descriptions.  More than
three quarters of these estimates are accurate;
every member of the family can be evaluated by a halting deterministic computation
using $O(B)$ space.  Exhaustive enumeration then gives a deterministic
answer.  The second guarantee is essential: enumeration must finish even on choices
for which the estimate is inaccurate.

\paragraph{The exact target.}
The configuration reduction of Section~\ref{sec:model} gives a strictly
forward substochastic matrix $S$ on a polynomial-size set $V_0$ and a
vector $e\in\{0,1\}^{V_0}$ indicating accepting final configurations.
The vector
\[
                         p_0=(I-S)^{-1}e
\]
records acceptance probabilities from all configurations.  The inverse
is a finite sum, because every transition advances time.
Section~\ref{sec:hierarchy} replaces $S$ successively by transitions
$C_l$ on sets $V_l$ of configuration copies, carrying accumulated
rewards $W_l\in[0,1]^{V_l}$.  Each copy step has a designated primary
copy.  If $z_l$ is obtained from the original start $x_0$ by always
choosing that copy, then
\[
       0\le p_0(x_0)-W_l(z_l)
                    \le |V_0|\,2^{-(H-1)l},
\]
where $H\ge4$ is fixed.  Thus $L=O(B)$ stages make $W_L(z_L)$ accurate
to any fixed inverse-polynomial error.  Every stage has at most
$N=2^{C_NB}$ vertices for a fixed constant $C_N$.

\paragraph{From the estimator to a deterministic decision.}
The algorithm estimates $W_L$ in one shared finite random environment
$\sigma$.  Its law is uniform on a specified finite set of matrix-and-bit-string
encodings of length $O(B)$.  On each channel a retention test
selects vertices; write $X_f(v)$ for the indicator that $v$ is retained
at the final stage rate $f=2^{-HL}$, and $\pi_f$ for its probability.
Section~\ref{sec:env} gives $\pi_f\asymp f$.
Theorem~\ref{thm:fp-estimators} proves, for the actual estimator $\bar W_L$,
\[
 \frac1N\mathbb E_\sigma\sum_{v\in V_L}\frac{X_f(v)}f
          |\bar W_L(v)-W_L(v)|^8\le C^8 2^{-8p}.
\]
Here $p$ is the statistical accuracy exponent.
Keeping only the summand at $z_L$ gives
\[
 \mathbb E\!\left[|\bar W_L(z_L)-W_L(z_L)|^8
                     \mid X_f(z_L)=1\right]
       \le N\frac f{\pi_f}C^8 2^{-8p}.
\]
The factor $N$ is the price of passing from the integrated bound to one
specified start.  Since $\log N=O(B)$, taking the estimator's statistical
budget to be a sufficiently large fixed multiple of $B$ makes $p$ large
enough to pay this price and gives an
accurate value on more than three quarters of the conditional
environments.

The implementation computes dyadic approximations to the same fixed
value $\bar W_L(z_L)$ on every environment.  Numerical precision
changes the requested digits, not the lists and ports defining that
value.  Choosing the stage count, statistical budget and digit accuracy
makes the total error less than $1/8$ on more than three quarters of
the encodings conditioned on $X_f(z_L)=1$.  Their median has the same
error bound.  We enumerate the original encodings with their original
multiplicities, so the multiset reproduces exactly this conditional law.

Repeated enumeration finds the median while retaining only the current
encoding and $O(B)$-bit numerical and counting registers.  Every digit
query used here requests one of the first $O(B)$ binary places,
terminates in $O(B)$ space, and confines
its physical input heads between the real endmarkers.  It therefore has only
$2^{O(B)}$ complete configurations and takes at most that many steps.
There are also $2^{O(B)}$ encodings, so the entire computation takes
polynomial time.  Comparing its answer with $1/2$ distinguishes
acceptance probabilities at most $1/3$ from those at least $2/3$.
This proves the nontrivial containment in Theorem~\ref{thm:main}.

\paragraph{Why the hierarchy can be estimated.}
The hierarchy reduces transition entries more rapidly than it increases
the number of vertices incident to positive entries.  A correction $E$ replaces $C$ by
$D=C-E+EC$, with
\[
                         I-D=(I+E)(I-C).
\]
This identity transfers the removed transition contribution into the
reward.  Detectors identify columns with many substantial entries and
suppress correction there.  The copy step distributes the remaining
large entries while controlling how many vertices are incident to them.  An ordered-support argument gives a factorial
truncation bound, so only constantly many correction steps are needed
at each stage.

The exact matrices and conditional expectations define the targets;
they are not stored.  Section~\ref{sec:atom} represents their
contributions by sample-dependent tables with bounded row support.
Corrections also have bounded incoming degree.  Their true conditional
means recover the required entries.  The algorithm substitutes earlier
estimates into these formulas, so the analysis must bound the resulting
bias without assuming that the estimated factors are independent.
Conditional averaging preserves the specified source and endpoint
retention tests.  Its uniform mixing bound is derived from classical
property~(T) results of Shalom~\cite{Shalom1999};
Section~\ref{sec:averaging} proves the finite-action and sampling
consequences used here.

At parent budget $M$, Section~\ref{sec:est} averages differences between
child estimates at consecutive budgets $m$ and $m-1$.  These differences
have small norm because both approximate the same true sample term.
Their contributions at the same endpoint are combined before measuring
the difference.
The slack $M-m$ has two roles: it determines the length of conditional
averaging available to reduce the noise, and it bounds the description
of an immediate sample path.  Grouping comparable slacks lets
Section~\ref{sec:buf-compression} replace each row array by a short list
of exceptions to a common default.  The proof charges deterministic
conditional-mean error separately from residual noise.

A product follows a correction edge to an intermediate vertex and then
queries a right factor there.  Releasing the old source changes the
weighted error domain.  Bounded incoming degree controls that change,
and the task's accuracy assignment pays the resulting factor.  Thus
the same construction that permits a source address to be discarded
also appears explicitly in the error analysis.

\paragraph{Why the lists fit in logarithmic space.}
A list entry is found by a path.  Retaining a full endpoint at every
recursive level would cost too much.  Section~\ref{sec:fp-blocks}
therefore compares endpoints by short fingerprint keys at small slacks,
combining polynomial
evaluation and affine hashing~\cite{Schwartz1980,CarterWegman1979}.
Candidate endpoints
can depend on earlier comparisons, so a fixed-pair collision estimate
alone is insufficient.  A comparison calculation using exact equality
within one budget block produces candidate sets independent of that
block's fingerprint bits.  Smooth support functions charge paths that
leave these sets to lower-call errors; collisions within the fixed sets
supply the remaining error.  This coupled induction proves the moment
bound for the actual estimator.

Sections~\ref{sec:arith} and~\ref{sec:ctrl} establish its total space
bound.  Numerical operations request one bounded digit at a time.
The controller keeps one vertex cursor, uses bounded inverse ports for
correction moves, and recovers incoming accesses by traversing finite
program trees, using the contour idea of Cook and
McKenzie~\cite{CookMcKenzie1987} and the space-preserving traversal
perspective of Lange, McKenzie and Tapp~\cite{LMT2000}.  A shared catalytic bit vector permits fingerprint
comparisons without keeping a key index at every recursive level.
The controller assigns each call a decreasing recursion allowance $w$,
with top allowance $O(B)$, and uses the requested digit place $t$ as its
numerical clock.  If the child has allowance $w'$ and digit place $j$,
all local data retained during that call occupy at most
\[
                         A(w-w')+D(t-j)
\]
bits, including the controller copies used by a traversal.  The child
clock may increase only in proportion to the decrease of allowance.
These differences telescope along the active call chain; the environment,
cursor, catalyst and scratch are shared once.  The bound and termination
do not depend on statistical accuracy.

Section~\ref{sec:conclusion} carries out the median construction for
fixed inverse-polynomial error and derives the class equality and
promise separation.  It then obtains one algorithm for all unary
precision requests by padding, and constructs certified accepting
computations when the machine's acceptance probability has an
inverse-polynomial lower bound.
Section~\ref{sec:compiler} is a separate application of fixed-precision
approximation: one resource-capped universal interpreter yields a fixed
deterministic library, whose syntactic specialization supplies explicit
resource bounds for each source machine.
